\documentclass[10pt,a4paper]{amsart}
\usepackage[margin=19mm]{geometry}
\usepackage{amsmath,amssymb,mathtools}
\usepackage[scr=rsfs,cal=euler]{mathalfa}
\usepackage{xfrac}
\usepackage[unicode,hidelinks,bookmarks=false]{hyperref}
\newcommand{\ie}{i.e.\ }
\newcommand{\cf}{cf.\ }
\newcommand{\resp}{respectively\ }
\newcommand{\Subsection}{\subsection}
\providecommand{\nobreakdash}{\nobreak\hbox{-}}
\setlength{\emergencystretch}{2em}
\multlinegap0pt
\usepackage{amssymb}
\usepackage{enumerate}
\usepackage[shortlabels]{enumitem}
\newtheorem{lemma}{Lemma}
\title{Diophantine Equations for Polynomial Recursive Sequences}
\author{Darsana N, Sudhansu Sekhar Rout}
\date{Source: arXiv:2512.20384v1 (CC BY 4.0)}
\begin{document}
\maketitle

\noindent\emph{Source and licence.} Diophantine Equations for Polynomial Recursive Sequences. Version arXiv:2512.20384v1 (CC BY 4.0), distributed by arXiv under the Creative Commons Attribution 4.0 International licence, http://creativecommons.org/licenses/by/4.0/, as recorded in the arXiv licence metadata for this exact version and shown on the abstract page https://arxiv.org/abs/2512.20384v1. Full licence notice: \texttt{licenses/arxiv-2512.20384-CC-BY-4.0.txt}.\par\smallskip

\noindent\textit{Editorial scope.} Counted unit: the article's lemma lemequality (source0.tex lines 424--430). The article's complete original proof of it is delivered with it (source0.tex lines 431--473, one \texttt{proof} environment of 43 lines). No mathematics is written, repaired or re-derived: the statement and the proof are the article's own lines, in the article's own notation, and no retained line is substituted or re-flowed.  No further source-local context is needed: every cross-reference of every retained line resolves inside the two delivered spans.  Nothing was omitted from any retained span, so the package carries no omission record.  No retained line contains a citation, so the wrapper carries no bibliography and no bibliography entry.  No retained line contains a figure environment, an image-inclusion command or drawing code, so no image is delivered and the package declares no asset dependency.  Editorial formatting only, in addition: an AMS \texttt{amsart} preamble that restates the article's own declarations and macros used by the retained lines, namely the environments \texttt{lemma} on separate counters, together with the standard packages those lines need.  The retained environments print in this document as Lemma 1 in order of appearance, whereas the article prints the same items with the numbering of its own full document.  The complete native source of the article as distributed in the arXiv e-print for this exact version is bundled unchanged as \texttt{sources/191560/source0.tex}.
    \begin{lemma}\label{lemequality}
    Let $f,g\in L^{*}$. Then
        \begin{itemize}
            \item[(i)] $\mathcal{H}(f+g)=\mathcal{H}(f)+\mathcal{H}(g)$, if $\nu(f)\neq \nu (g)$ and $\max \{\nu(f),\nu (g)\}>0$ for all valuations $\nu$ over $L$.
            \item[(ii)] $\mathcal{H}(fg)=\mathcal{H}(f)+\mathcal{H}(g)$, if $\nu(f)\nu (g)>0$ for all valuations $\nu$ over $L$. 
        \end{itemize}
    \end{lemma}

    \begin{proof}
        \begin{enumerate}
            \item[($i$)] Recall that
            \[\mathcal{H}(f+g)=- \sum_{\nu} \min \{ 0,\nu (f+g)\}.\]
            Therefore, it suffices to verify that for every valuation $\nu$, 
            \begin{equation}\label{eqnew1}
                \min \{ 0,\nu (f+g)\}=\min \{ 0,\nu (f)\}+\min \{ 0,\nu (g)\}.
            \end{equation} Since 
            \begin{equation*}
                \nu(f+g)\geq \min\{\nu(f),\nu(g)\}.
            \end{equation*}and because we assume $\nu(f)\neq \nu(g)$, we infact have 
            \begin{equation*}
                \nu(f+g)= \min\{\nu(f),\nu(g)\}.
            \end{equation*} Without loss of generality assume that $\nu(f)<\nu(g)$; hence $\nu(f+g)=\nu(f)$. Fix a valuation $\nu$. Now we consider the following three cases.\\
            \textit{Case 1:} Suppose  $\nu(f+g)=0.$ Then $\nu(f)=0$ and $\nu(g)>0$. Thus
            \[
            \min \{0, \nu(f+g)\}=0=\min \{0, \nu(f)\}+\min \{0, \nu(g)\}.
            \]  So \eqref{eqnew1} holds.\\
            \textit{Case 2:} Suppose $\nu(f+g)>0$. Then $\nu(f)>0$ and since $\nu(f)<\nu(g)$, also $\nu(g)>0$. Hence, in \eqref{eqnew1} $\nu(f+g),\nu(f),\nu(g)$ are positive, and therefore is zero and thus equality holds.\\
            \textit{Case 3:} Suppose $\nu(f+g)<0$. Then $\nu(f)<0$. Also note that either $\nu(g)<0$ or $\nu(g)>0$. Since by our assumption $\max \{\nu(f), \nu(g)\}=\nu(g)>0$ and therefore $\min \{0,\nu (g)\}=0$. Thus,
            \[
            \min\{0,\nu(f+g)\}=\nu(f)=\min\{0,\nu(f)\}+\min\{0,\nu(g)\}
            \]and\eqref{eqnew1} holds. This completes the proof of part $(i)$.
            \item[($ii$)] We have $\nu(fg)=\nu(f)+\nu(g)$. Therefore,
            \begin{equation*}
                \mathcal{H}(fg)=-\sum_{\nu} \min \{0,\nu(fg)\}= -\sum_{\nu} \min \{0,\nu(f)+\nu(g)\}. 
            \end{equation*} So it is suffices to verify that for every valuation $\nu$,
            \begin{equation}\label{eqnew2}
                \min \{ 0,\nu (f)+\nu(g)\}=\min \{ 0,\nu (f)\}+\min \{ 0,\nu (g)\}.
            \end{equation} We check this by considering all possible sign combination of $\nu(f)$ and $\nu(g)$.\\
            \textit{Case 1:} Suppose $\nu(f)=\nu(g)=0$. Then, $\nu(f)+\nu(g)=0$, and the identity \eqref{eqnew2} holds trivially. \\
            \textit{Case 2:} Suppose $\nu(f)>0$ and $\nu(g)=0$. Then $\nu(f)+\nu(g)=\nu(f)$ and hence all the minimum in $\eqref{eqnew2}$ are zero and the identity holds.\\
            \textit{Case 3:} Suppose $\nu(f)>0$ and $\nu(g)>0$. Then 
            \[
            \min \{0, \nu(f+g)\}=0=\min \{0, \nu(f)\}+\min \{0, \nu(g)\}.
            \]
            \textit{Case 4:} Suppose $\nu(f)<0$ and $\nu(g)<0$. Then $\nu(f)+\nu(g)<0$. Thus, 
            \begin{equation*}
                \min \{ 0,\nu(f)+\nu(g)\}= \nu(f)+\nu(g)= \min \{ 0,\nu(f)\}+\min \{0,\nu(g)\}.
            \end{equation*}
             Similarly, if we interchange the roles of $f$ and $g$ in all the above cases, we get the same conclusion. This completes the proof of Lemma \ref{lemequality}.
        \end{enumerate} 
    \end{proof}

\end{document}
